L'espace est muni d'un repère orthonormé \Oijk.

On considère :

\begin{itemize}[label=$\bullet~$]
\item les points A(4~;~2~;~2), B$(5~;~-2~;~3)$ et C(1~;~1~;~1) ;
\item la droite $\Delta$ dont une représentation paramétrique est donnée par

\[\left\{\begin{array}{lc l}
x	&=& 1 + 2t\\
y 	&=&1+ \phantom{2}t\\
z 	&=& 1 +2t\\
\end{array}\right. \text{avec}\: t \in \R\]

\item le plan $\mathcal{P}$ contenant le point A et perpendiculaire à la droite $\Delta$.
\end{itemize}

\medskip

\begin{enumerate}
\item Vérifier que la droite $\Delta$ contient le point C(1~;~1~;~1) mais pas le point A.
\item
	\begin{enumerate}
		\item Démontrer qu'une équation cartésienne du plan ~ est $2x + y + 2z - 14 = 0$.
		\item Vérifier que le plan $\mathcal{P}$ contient le point B mais pas le point C.
	\end{enumerate}
\item On considère le point D(3~;~2~;~3).
	\begin{enumerate}
		\item Démontrer que le point D est le projeté orthogonal du point C sur le plan $\mathcal{P}$.
		\item Justifier que les points A, B, C et D ne sont pas coplanaires.
		\item Calculer le produit scalaire $\vect{\text{AB}} \cdot \vect{\text{AD}}$.
		\item Calculer le volume du tétraèdre ABCD.
		
On rappelle que le volume d'un tétraèdre est donné par $V = \dfrac13 \times B \times h$ où $B$ est l'aire d'une base du tétraèdre et $h$ la hauteur relative à cette base.

	\end{enumerate}
\item On appelle H le projeté orthogonal du point A sur la droite (BC).
	\begin{enumerate}
		\item Vérifier que les coordonnées du point H sont $\left(\dfrac{73}{29}~;~\dfrac{-4}{29}~;~\dfrac{51}{29}\right)$.
		\item Démontrer que l'aire du triangle ABC est $\dfrac{3\sqrt{22}}{2}$.
		\item En déduire la distance du point D au plan (ABC).
	\end{enumerate}
\end{enumerate}

\medskip
